\section{Background and Related Work}
\label{sec:related}

Prior work touches co-evolutionary security from several directions, but
rarely combines all of the following: autonomous agents on both sides, a
cybersecurity environment, repeated episodes, explicit adaptation levels,
and reproducible measurement of strategy change. We review five strands and
 summarise the gap in Table~\ref{tab:gap}.

\paragraph{LLM agents and tool use.}
ReAct-style loops couple reasoning with acting through external
tools~\cite{yao2023react,schick2023toolformer}. Multi-agent systems such as
AutoGen and generative-agent societies show that language-model instances
can plan, delegate and persist memory~\cite{wu2023autogen,park2023generative,xi2023rise,shinn2023reflexion}.
These architectures increase the attack surface (tools, memory, messages)
but do not, by themselves, define a co-evolutionary evaluation protocol.
We adopt a deliberately simple decision loop so that differences across
runs can be traced to adaptation level rather than planner design. Planner
complexity (multi-agent debate, retrieval, code synthesis) is orthogonal
to the question we ask: whether \emph{memory regime} and \emph{opponent
adaptation} change measurable security outcomes when the tool catalogue
and detection physics are held fixed.

\paragraph{Security of agentic systems.}
Indirect prompt injection~\cite{greshake2023not}, jailbreaks~\cite{zou2023universal}
and systematic red teaming~\cite{perez2022red} show that models can be
steered into unintended actions. Agent-specific benchmarks (AgentDojo,
InjecAgent, AgentPoison) extend this to tool misuse, memory poisoning and
environment injection~\cite{debenedetti2024agentdojo,zhan2024injecagent,chen2024agentpoison}.
Multi-agent studies document collusion and steganographic
coordination~\cite{abdelnabi2024cooperation}. The typical evaluation
freezes the defender or truncates the horizon; it does not measure what
happens when a defensive agent also remembers and adapts.

\paragraph{Autonomous cyber defence and cyber ranges.}
CybORG/CAGE, CyberBattleSim and related gyms let learning agents operate
on emulated networks~\cite{standen2021cyborg,microsoft2021cyberbattlesim,molina2021network}.
MITRE ATT\&CK supplies vocabulary for adversary
behaviour~\cite{strom2018attck}. These platforms emphasise high-fidelity
environments and defender training against scripted or slowly varying
attackers. Co-evolution of \emph{both} sides under declared memory regimes,
with explicit emergence criteria, is outside their usual experimental
design. Our simulator backend is intentionally smaller in topology so that
factorial sweeps over adaptation remain tractable; the optional Kubernetes
path preserves interface compatibility for later fidelity studies.

\paragraph{Security games and multi-agent learning.}
Stackelberg security games allocate defensive resources against a best
responding attacker~\cite{tambe2011security}. Multi-agent reinforcement
learning studies cooperation and competition under partial
observability~\cite{lowe2017multi,foerster2018counterfactual,kaelbling1998planning}.
Our setting is closer to the latter (both sides learn; no closed-form best
response), but grounded in a cyber environment with structured tools and
logged trajectories rather than an abstract matrix game. Unlike classical
repeated games with known payoffs, tool outcomes here depend on hidden host
state and stochastic detection, which motivates reporting DR and CEP alongside
any success metric.

\paragraph{Evolutionary computation.}
Co-evolution as a machine-learning method dates to competitive co-evolved
sorting networks and fitness sharing~\cite{hillis1990co,rosin1997new};
evolutionary game theory provides replicator dynamics and evolutionarily
stable strategies~\cite{hofbauer2003evolutionary,weibull1995evolutionary}.
We do not evolve network weights. We study \emph{behavioural}
co-evolution: policies change because agents remember outcomes and alter
subsequent tool choices. The analogy is a security arms race expressed
through strategy signatures and episode dynamics, not neuroevolution.

\begin{table}[!t]
\caption{Related directions and the limitation this work addresses.}
\label{tab:gap}
\centering
\begin{tabularx}{\columnwidth}{@{}lX@{}}
\toprule
\textbf{Direction} & \textbf{Typical limitation} \\
\midrule
Single-agent security~\cite{greshake2023not,debenedetti2024agentdojo}
  & No strategic interaction \\
Agent-to-agent attacks~\cite{zhan2024injecagent}
  & Relatively static adversary \\
Multi-agent collusion~\cite{abdelnabi2024cooperation}
  & Limited long-horizon adaptation \\
Autonomous cyber defence~\cite{standen2021cyborg}
  & Defender-centric scoring \\
Security games~\cite{tambe2011security}
  & Closed-form or scripted best response \\
Evolutionary AI~\cite{rosin1997new}
  & Rarely a cyber environment \\
\midrule
This work & Both sides adapt; cyber env.; measured co-evolution \\
\bottomrule
\end{tabularx}
\end{table}

The gap is therefore not a missing benchmark feature but a missing
\emph{experimental question}: what changes when adaptation is a first-class
parameter on both sides of a repeated cyber interaction?
